% ID: SEL-D0079-OB-E03
% ESTADO: cotejada. Segunda pasada visual y borrador compilado revisados.
% Fuente DOC: Selectividad/sel_2012_matematicas_aplicadas/1 MATEMÁT. APLIC. CC. SS/EXAMEN 1-ANDALUCÍA.doc
% SHA-256 original: a5dfff58dac253bf021b0d2b9409d68987e9ce6e97cfa4abbaf840019a7a18e2
% Página convertida: 2; opción B; ejercicio 3.

Lanzamos un dado, si sale 5 o 6 extraemos una bola de una urna A, que contiene 6 bolas blancas y 4 negras. Si sale otro resultado se extrae una bola de la urna B, que contiene 3 bolas blancas y 7 negras. Calcule:
\begin{enumerate}[label=\alph*)]
\item \textbf{(1 punto)} La probabilidad de que la bola extraída sea negra.
\item \textbf{(0.5 puntos)} La probabilidad de que la bola sea negra y de la urna B.
\item \textbf{(1 punto)} La probabilidad de que haya salido menos de 5 si la bola extraída ha sido blanca.
\end{enumerate}
